
Four experiments, each addressing a specific research question:

\textbf{RQ1} (Cross-architecture SSL transfer): Does graph-based SSL with permutation equivariance outperform flat tokenization (SANE) on held-out ViT model accuracy prediction?

\textbf{RQ2} (Distribution shift measurement): Does graph encoding reduce measurable distribution shift between CNN training zoo and ViT test zoo latent representations?

\textbf{RQ3} (Symmetry group ablation): Does scale+permutation equivariance improve over permutation-only equivariance for ViT cross-architecture transfer?

\textbf{RQ4} (Latent interpolation): Does latent-space model interpolation produce more accurate models than weight-space averaging?

\subsubsection{Datasets}

\textbf{Training zoo.} Approximately 3,000 CNN checkpoints from the SANE MultiZoo CIFAR-10 subset (tune\_zoo\_cifar10\_uniform\_small) [Schurholt2024Towards]. No ViT models appear during encoder training. The full SANE MultiZoo contains approximately 30,000 CNN checkpoints; we use a 3,000-model subset due to computational constraints.

\textbf{Test zoo.} 53 ViT-S/16 ImageNet checkpoints from the ViT Model Zoo [Falk2025ModelZoo]. The original dataset contains 250 models; we use 53 due to local availability constraints. All reported cross-architecture $R^2$ values are computed on these 53 models.

\textbf{Interpolation dataset (RQ4).} 200 real CNN checkpoints from the SANE ModelZoo (zenodo:13144018), forming 501 model pairs.

\subsubsection{Baselines}

\begin{itemize}
\item \textbf{SANE} [Schurholt2024Towards]: Flat tokenization + VICReg fix.
\item \textbf{EquiSSL} (scale+permutation, ScaleGMN [Kalogeropoulos2024Scale]): Ablation variant with full monomial matrix group equivariance. Trained for 50 epochs (h-e1 checkpoint; see Section 4 note below).
\item \textbf{EquiSSL-perm} (permutation-only, neural-graphs [Kofinas2024Graph]): Primary method. Trained for 100 epochs (h-m1 checkpoint).
\item \textbf{Weight-space averaging} (RQ4): Naive model merging baseline --- ($\theta _A + \theta _{B})/2$ --- evaluated on CIFAR-10.
\end{itemize}

\textbf{Note on training duration discrepancy.} EquiSSL (ScaleGMN) was trained for 50 epochs in the h-e1 experiment; EquiSSL-perm was trained for 100 epochs in h-m1. The h-m2 ablation experiment (Section 5.2, Table 2) reuses these pre-existing checkpoints rather than training new encoders. This creates a training duration confound in Table 2: the $\Delta R^{2}=+0.143$ advantage for EquiSSL-perm in h-m2 may partly reflect the additional 50 training epochs rather than the symmetry group difference alone. The h-m1 comparison (Table 1), where both encoders were evaluated under the same probe setup, shows a smaller directionally consistent advantage ($\Delta R^{2}=+0.021$ for EquiSSL-perm). A fully controlled ablation would train both models for identical epoch counts.

\subsubsection{Evaluation Metrics}

\textbf{$R^2$}: Primary metric for property prediction (RQ1--RQ3). \textbf{MMD} (RBF kernel with median heuristic): Secondary metric for distribution shift (RQ2). \textbf{Accuracy delta}: mean(acc\_latent $-$ acc\_ws) over 501 pairs (RQ4).

